% !PRINTED 196
\weberpage{196}
That the signs of $\sqrt{-3\,D}$ are to be taken, with regard to the signs in (11), as has been done in the formulas (13), follows most simply if one carries out the computation for the special case $a_0 = 0$, $a_3 = 0$.

\begin{center}
\S.\,63.
\end{center}

\begin{center}
\gesperrt{The full system of forms of the binary cubic form.}
\end{center}

We wish to show further that the invariants and covariants of the cubic form are exhausted by the functions $D, f, H, Q$, i.\,e. that all invariants and covariants of a cubic form can be expressed rationally in terms of these four forms.

So let
\[
C = C(x,y,a)
\]
be a covariant of degree $\nu$ in the variables and of degree $\mu$ in the coefficients, of which we shall assume that it has only integer numerical coefficients.

Then, for any transformed form,
\weq{1}{C' = C(\xi,\eta,a') = r^{\lambda}\,C(x,y,a),}
\weq{2}{\lambda = \frac{3\,\mu - \nu}{2}.}

We now understand by $\xi$, $\eta$ the variables of the normal form considered above. If $\varepsilon$ denotes a third root of unity, then
\[
\xi = \varepsilon\,\xi',\quad \eta = \varepsilon^2\,\eta'
\]
is a linear substitution with determinant 1, by which the normal form \S.\,62, (9) passes into itself. It follows, therefore, that $C'$ cannot change when $\xi$, $\eta$ are replaced by $\varepsilon\,\xi$, $\varepsilon^2\,\eta$; hence $C'$ can depend rationally only on $\xi\,\eta$, $\xi^3$, $\eta^3$.

The permutation of $\xi$ with $\eta$ corresponds to a substitution with determinant $-1$; under it, therefore, $C'$ remains unchanged or changes its sign, according as $\lambda$ is even or odd. Hence $C'$ consists of a sum of terms of the form
\[
M\,(\xi\,\eta)^h\,(\xi^{3k} \pm \eta^{3k}),
\]
in which $M$ is an integer factor (since we assume the numerical coefficients in $C$, and hence also in $C'$, to be
%CONT>>
